\section{结果}
\label{sec:results}
\subsection{工程门槛与正式数据覆盖}
预试验中，M0 的三个 aBN1 响应依次为 10、10、10，2 s 休息后为 10；M1-zero 相同，且两者的完整全脑训练脉冲表一致。M1 的对应响应为 8、5、3，休息后为 6。三个模型各自保存并恢复训练末态后，恢复分支的完整脉冲表均能一致重放。预试验表明实现能够保留状态，并通过“零强度扩展不改变结果”的工程门槛；它没有被当作独立生物验证集。

正式实验完整覆盖 20 条序列，含 240 个训练响应和 40 个恢复响应，共 280 个窗口。重新回读 60 个原始事件文件后，各读出计数、总脉冲数及感觉脉冲数均与登记结果一致。同一种子的感觉神经元相对发放时刻，在模型、间隔、重复刺激和恢复探测之间完全相同。登记文件中的窗口外脉冲数为零；这不是程序把状态清零的结果，连续状态仍由方程推进。

\subsection{M0 在两个间隔下均未产生逐次递减}
图\ref{fig:training}与表\ref{tab:main}显示，固定连接模型的 aBN1 响应逐次恒定。种子 1701--1704 的每次计数均为 10，种子 1705 均为 9；两种间隔下组均值均为 9.80，早期和末期相等，所有 $D=0$。2 s、10 s 探测同样保持各自的初始计数。

aDN1 与 aDN2 也在各序列内保持恒定。这个结果不是“不够显著的下降”，而是在当前确定性回放输入和所报告计数精度下，根本没有观察到跨次计数变化。它支持一个有限结论：\textbf{在所测试的 JO-CE 输入、原始参数和两个间隔下，M0 没有表现出本研究定义的响应递减。}

\newcommand{\TrainingPanel}[2]{%
\begin{tikzpicture}
\begin{axis}[width=\linewidth,height=5.6cm,title={#2},
 xlabel={刺激序号},ylabel={aBN1 脉冲数},xmin=0.7,xmax=12.3,ymin=-0.3,ymax=11,
 xtick={1,3,6,9,12},ytick={0,2,4,6,8,10},grid=major,grid style={black!10},
 tick label style={font=\footnotesize},label style={font=\small},title style={font=\small\bfseries}]
\addplot[name path=staticlow,draw=none,forget plot] table[x=trial,y=low]{data/train_static_#1.csv};
\addplot[name path=statichigh,draw=none,forget plot] table[x=trial,y=high]{data/train_static_#1.csv};
\addplot[Gray!15,forget plot] fill between[of=staticlow and statichigh];
\addplot[name path=stdlow,draw=none,forget plot] table[x=trial,y=low]{data/train_std_#1.csv};
\addplot[name path=stdhigh,draw=none,forget plot] table[x=trial,y=high]{data/train_std_#1.csv};
\addplot[Blue!15,forget plot] fill between[of=stdlow and stdhigh];
\addplot[Gray,thick,mark=square*,mark size=1.5pt] table[x=trial,y=mean]{data/train_static_#1.csv};
\addplot[Blue,thick,mark=*,mark size=1.6pt] table[x=trial,y=mean]{data/train_std_#1.csv};
\end{axis}
\end{tikzpicture}}
\begin{figure}[tbp]
\centering
\begin{minipage}{0.49\linewidth}\TrainingPanel{500}{A. 起始间隔 0.5 s}\end{minipage}\hfill
\begin{minipage}{0.49\linewidth}\TrainingPanel{2000}{B. 起始间隔 2 s}\end{minipage}
\par\smallskip{\small\textcolor{Gray}{$\blacksquare$ M0：固定连接}\qquad\textcolor{Blue}{$\bullet$ M1：局部 STD}}
\caption{重复刺激下 aBN1 的响应。每点为五个冻结输入实现的均值；阴影为最小值至最大值，不是置信区间。计数窗口为每次刺激起始后的 300 ms。两幅图纵轴相同。M0 逐次恒定；M1 在短间隔下强烈递减，在长间隔下保留较大响应。横轴是刺激序号，不是相同长度的物理时间。}
\label{fig:training}
\end{figure}


\begin{table}[tbp]
\centering\small
\setlength{\tabcolsep}{5pt}
\caption{主要读出 aBN1 的结果。除 $D$ 外单位均为每个 300 ms 窗口的脉冲数；每格为五个输入实现的均值。$I,E,L$ 分别为首次、前 3 次和后 3 次响应，$P_t$ 为休息 $t$ 后的独立探测。$D$ 先按每个种子计算再平均。}
\label{tab:main}
\begin{tabular}{lrrrrrr}
\toprule
模型 / 间隔 & $I$ & $E$ & $L$ & $D$（\%） & $P_2$ & $P_{10}$\\
\midrule
\inputrows{data/main_rows.tex}
\bottomrule
\end{tabular}
\end{table}

\subsection{M1 在短间隔下产生强递减与休息恢复}
在 $T=0.5\secunit$ 条件下，M1 首次响应均值为 8.20，前 3 次均值为 5.40，后 3 次均值降至 0.53。逐种子递减比例均值为 90.22\%，范围为 81.25\%--100\%；5/5 输入实现通过 H1 筛查，也全部通过平台筛查。

短间隔组 2 s 探测均值为 5.40，10 s 探测均值为 8.20（图\ref{fig:recovery}）。所有五例都满足预设 H2 门槛。但组均值恢复不意味着每个实现逐点完全恢复：10 s 探测相对首次的比值为 87.5\%--112.5\%。这些略低或略高于首次的计数不能被删去或截断为“100\%”；本文保留其原值，不用单一资源指数直接解释所有网络输出细节。

还应注意，M1 首次响应已经低于 M0 的 9.80，均值降幅约为 16.3\%。因此，该机制不仅储存跨刺激历史，也影响第一段 200 ms 刺激内部的传递。若只展示每组各自归一化后的下降曲线，就可能掩盖这项正常功能损失。

\begin{primer}{90\% 下降到底相对什么}
这里的 90.22\% 是“后 3 次平均响应，相对前 3 次平均响应”的下降；不是相对未加 STD 的 M0，也不是相对第一次。第一例短间隔的 $E=(8+5+3)/3=5.33$、$L=1$，所以 $D=1-1/5.33=81.25\%$。先算每例再平均，得到组结果。这个明确的分母决定了百分数的意义。
\end{primer}

\newcommand{\RecoveryPanel}[2]{%
\begin{tikzpicture}
\begin{axis}[width=\linewidth,height=5.7cm,title={#2},ylabel={aBN1 脉冲数},xmin=-0.4,xmax=3.4,ymin=-0.3,ymax=11,
 xtick={0,1,2,3},xticklabels={首次,末期,休息2秒,休息10秒},ytick={0,2,4,6,8,10},
 grid=major,grid style={black!10},tick label style={font=\scriptsize},label style={font=\small},title style={font=\small\bfseries}]
\addplot[Gray,only marks,mark=square*,mark size=2pt,error bars/.cd,y dir=both,y explicit]
 table[x expr=\thisrow{stage}-0.07,y=mean,y error minus=minus,y error plus=plus]{data/recovery_static_#1.csv};
\addplot[Blue,only marks,mark=*,mark size=2pt,error bars/.cd,y dir=both,y explicit]
 table[x expr=\thisrow{stage}+0.07,y=mean,y error minus=minus,y error plus=plus]{data/recovery_std_#1.csv};
\end{axis}
\end{tikzpicture}}
\begin{figure}[tbp]
\centering
\begin{minipage}{0.49\linewidth}\RecoveryPanel{500}{A. 训练间隔 0.5 s}\end{minipage}\hfill
\begin{minipage}{0.49\linewidth}\RecoveryPanel{2000}{B. 训练间隔 2 s}\end{minipage}
\par\smallskip{\small\textcolor{Gray}{$\blacksquare$ M0：固定连接}\qquad\textcolor{Blue}{$\bullet$ M1：局部 STD}}
\caption{首次、末期及独立恢复探测的响应。点为五个输入实现的均值，误差棒为 min--max；“末期”先在每个实现内平均第 10--12 次，再跨实现汇总。横轴是四个离散条件，\textbf{不是均匀时间轴}，因此不连接为恢复曲线。10 s 时 M1 组均值回到首次水平，但个别实现低于或高于首次。两个恢复点不足以估计完整恢复动力学。}
\label{fig:recovery}
\end{figure}


\subsection{刺激间隔改变效应强度，但尚不能确定恢复速率规律}
在 $T=2\secunit$ 条件下，M1 的首次均值同样为 8.20，早期均值为 7.33，末期均值为 6.60；平均 $D$ 只有 9.88\%，范围为 0--18.18\%。五例均未达到 30\% 的 H1 门槛，4/5 通过平台筛查。较频繁刺激引起更强递减的方向在五个配对输入实现中一致（图\ref{fig:paired}），支持本范围内的 H4(a) 描述。

长间隔组的 2 s、10 s 探测均值分别为 6.40、8.20。不能因为 10 s 时两组都接近初始均值，就声称恢复时间常数相同；同样不能因为短间隔恢复的绝对增量更大，就声称其恢复更快。两组训练末态的损失深度不同，长间隔还有一例未通过平台标准，而且仅有两个恢复点。

经典 H4(b) 关心达到渐近水平后的频率依赖恢复。当前数据没有建立这一特征，也不足以严格否定它。M1 中所有局部资源确实共享同一个规定的恢复时间常数，但神经输出还要经过循环网络和阈值非线性；资源恢复速率不能不经检验地替代神经输出恢复速率。

\begin{figure}[tbp]
\centering
\begin{tikzpicture}
\begin{axis}[width=10.2cm,height=5.8cm,xmin=0.7,xmax=2.3,ymin=-4,ymax=106,
 xtick={1,2},xticklabels={0.5 s 间隔,2 s 间隔},ytick={0,20,40,60,80,100},ylabel={递减比例 $D$（\%）},
 grid=major,grid style={black!10},tick label style={font=\small},label style={font=\small}]
\pgfplotstableread{data/paired_decrement.csv}\pairedtable
\pgfplotsinvokeforeach{0,1,2,3,4}{%
 \pgfplotstablegetelem{#1}{D500}\of\pairedtable\edef\firstvalue{\pgfplotsretval}
 \pgfplotstablegetelem{#1}{D2000}\of\pairedtable\edef\secondvalue{\pgfplotsretval}
 \addplot[Blue!70,thick,mark=*,mark size=2pt] coordinates {(1,\firstvalue) (2,\secondvalue)};
}
\addplot[Gray,dashed,thick,mark=square*,mark size=2pt] coordinates {(1,0) (2,0)};
\addplot[Orange,densely dotted,forget plot] coordinates {(0.75,30) (2.25,30)};
\node[anchor=west,font=\footnotesize,text=Orange] at (axis cs:0.75,34) {预设 30\% 筛查阈值};
\end{axis}
\end{tikzpicture}
\caption{同一种子内的配对间隔比较。五条蓝线分别对应 M1 的五个输入实现；灰色虚线表示所有 M0 实现的 $D=0$。每条线连接同一输入在两种刺激间隔下的描述性效应量，不表示一个连续的参数拟合曲线。短间隔的递减在五个实现中都大于长间隔。部分端点重叠；具体数值见附录表\ref{tab:individual}。}
\label{fig:paired}
\end{figure}


\subsection{下游读出的地板效应限制了行为解释}
表\ref{tab:downstream}显示，M1 不仅改变 aBN1，也降低了相关下行输出。短间隔末期的 aDN1 和 aDN2 平均计数都降至零。但 aDN1 的 M1 首次平均已经只有 1.20，范围 0--2，其中一个输入实现首次就没有脉冲。

这造成重要解释限制：一个原本接近零的读出，很难用相对下降推断“保留正常行为能力的选择性习惯化”。本研究没有建立脉冲计数到梳理发生概率、动作幅度或身体运动的映射，也没有直接输出刺激来检验执行能力。因此，主要结论应停留在 aBN1 的习惯化样神经响应，而不应写成“数字果蝇已经学会忽略触碰”。

\begin{table}[tbp]
\centering\small
\caption{下游读出暴露正常传递与地板问题。首次为第 1 次响应；末期先平均每例第 10--12 次，再跨五个输入实现平均。单位为脉冲数。}
\label{tab:downstream}
\begin{tabular}{lrrrr}
\toprule
模型 / 间隔 & aDN1 首次 & aDN1 末期 & aDN2 首次 & aDN2 末期\\
\midrule
\inputrows{data/downstream_rows.tex}
\bottomrule
\end{tabular}
\end{table}

\subsection{实际覆盖的习惯化特征}
\begin{table}[tbp]
\centering\small
\caption{习惯化特征与本轮证据。H 编号沿用经典综述的特征顺序\citep{rankin2009}；筛查阳性不等价于生物行为确证。}
\label{tab:hallmarks}
\begin{tabularx}{\linewidth}{p{3.7cm}Y}
\toprule
特征 & 当前证据状态\\
\midrule
H1 反应递减 & M1 短间隔 aBN1 达到预设门槛；M0 阴性；M1 长间隔效应较弱。\\
H2 自发恢复 & M1 短间隔通过神经输出筛查；使用独立恢复分支。\\
H3 跨轮次增强 & 未做多轮训练--恢复循环，未检验。\\
H4(a) 频率依赖递减 & 两间隔配对比较支持更频繁刺激导致更强递减。\\
H4(b) 频率依赖恢复 & 未建立；时间点稀疏、末态深度不同、平台条件不完全满足。\\
H5 强度依赖 & 仅测试一个输入率和脉冲宽度，未检验。\\
H6 渐近后的继续积累 & 未做专门的渐近后追加刺激对照。\\
H7 特异性/泛化 & 没有加入可比较的新刺激 B，未检验。\\
H8 去习惯化 & 没有 A--B--A 和相应时间/敏化对照，未检验。\\
H9 去习惯化本身递减 & 未检验。\\
H10 长期习惯化 & 未引入长期机制，也未进行长期协议。\\
\bottomrule
\end{tabularx}
\end{table}

\subsection{执行中断、复核和资源}
首次执行达到 1,800 s 工具时限后停止。其 manifest 保留 \texttt{interrupted} 状态，不把该批次伪装为完整成功。已完成的 M0 十条序列和 M1/1701 两条序列被保留；尚未完成的 M1/1702/T500 片段未进入汇总。随后两批按原参数补完剩余八条序列，合并覆盖检查无缺失、无重复。

有效序列所在检查点及批次的记录耗时合计约 3,065 s（51.1 min），包含加载、编译、快照、静默间隔计算和写盘，且不包含中断后未完成部分。因此它不是纯积分吞吐率，也不是所有实际占用时间的精确总和。两次接续记录的峰值 RSS 约 2.47 GiB；第一次中断未保存最终峰值，不能用接续值冒充所有批次的全局上限。

在整理本报告时，再次从原始事件运行独立分析器，重建的报告对象与原有结果完全一致。数据制表程序又将主表数值从逐种子指标重新计算，并逐项核对；没有重新模拟、拟合或挑选种子来改善论文图形。
\FloatBarrier
